% \documentclass[11pt]{article}

% %%%% CUSTOM PREAMBLE
% \usepackage{./preamble}

% %% ADD MISC DEFINITIONS HERE %%%%%%%%%%%%%%%%%%%%%%%%%%%%%%%%%%%

% %%%%%%%%%%%%%%%%%%%%%%%%%%%%%%%%%%%%%%%%%%%%%%%%%%%%%
% %\addtolength{\topmargin}{-1cm}
% %\addtolength{\textheight}{2cm}

% \begin{document}

%\maketitle

\newpage
 %%*****************************************************
\noindent
\textbf{CS4268 AY2025-26 Sem 2}
\\
\noindent
Lecturer: Divesh Aggarwal
\hfill
Lecture 4                      %%% FILL IN LECTURE NUMBER HERE
\hfill
Date: February 5, 2026          %%% FILL IN LECTURE DATE HERE


\noindent
\rule{\textwidth}{1pt}

\medskip

%%%%%%%%%%%%%%%%%%%%%%%%%%%%%%%%%%%%%%%%%%%%%%%%%%%%%%%%%%%%%%%%
%% BODY OF SCRIBE NOTES GOES HERE
%%%%%%%%%%%%%%%%%%%%%%%%%%%%%%%%%%%%%%%%%%%%%%%%%%%%%%%%%%%%%%%%

%%%Write lecture topic above.
\textbf{Instructions}\\
\begin{itemize}
    \item Try to leave the preamble unchanged, only make changes to \texttt{main.tex} and  \texttt{references.bib}. Add images to the \texttt{images} folder. 
    \item Share one (zipped) directory preserving the directory structure on Canvas. 
    \item Prepare organized notes by elaborating on the broad items provided in different sections below based on the lecture material completing any missing details from lectures.
    \item Feel free to modify structure and/or order of contents as you see fit; the below is just a guideline.
\end{itemize} 

\section{The No-Cloning Theorem and Quantum Teleportation}

\subsection{No Quantum Clones}
\kk{add full details}


\subsection{Quantum Teleportation}

Quantum Teleportation
Idea: Alice has a state $\ket{\psi}$ (whose amplitudes she doesn't even know), she sends it to Bob using a shared $\ket{EPR}$ and two bits of classical communication.

Method: Starting from $\ket{\psi}\ket{EPR}$, Alice applies $H\circ CNOT$ on her two qubits to obtain
\begin{align*}
    \frac{1}{2}[&\ket{00}(\alpha\ket{0}+\beta\ket{1}) + \ket{01}(\alpha\ket{1}+\beta\ket{0})\\
    + &\ket{10}(\alpha\ket{0}-\beta\ket{1}) + \ket{11}(\alpha\ket{1}-\beta\ket{0})].
\end{align*}
Then measure her two qubits and send her results to Bob, who acts accordingly (apply $Z^{M_1}X^{M_2}$ on his qubit) to recover $\ket{\psi}$.

\kk{please add any further details you deem helpful, and the relevant circuits for this task}



\section[\texorpdfstring{Classical Computation $\bm{\subseteq}$ Quantum Computation}{}]{Classical Computation $\bm{\subseteq}$ Quantum Computation} 

\subsection{Circuit Model of Computation}

Classical circuits as model of computation (equivalent to turing machines model)

Add universal set of circuit gates, Shannon's universality theorem

most boolean functions are exponentially complex, add details and proof (if divesh covers). proof is combinatorial, essentially consider all possible ways to inject layers of gates, then show this doesnt cover that many functions if only a few layers.


\subsection{Reversible Computation}

Describe how to implement the basic gates like AND NOT OR reversibly all in terms of toffoli gate. 

Toffoli is unitary! thus classical computing $\subseteq$ quantum computing

Statement:

Given a function $f: \{0,1\}^n \rightarrow \{0,1\}^m$. Suppose we have a classical Boolean circuit $C_f$ implementing $f$ which requires $k$ elementary gates. We can convert these gates into an $O(k)$-sized reversible circuit $R_f$ comprising NOT and Toffoli gates, this also uses $O(k)$ ancillas. By copying $f(x)$ over to $y$ using CNOTs, then uncomputing, we get the reversible circuit:
\begin{align*}
    O_f\ket{x}\ket{y} = \ket{x}\ket{y \oplus f(x)}
\end{align*}
where we have omitted the ancillas. Roughly speaking, $O_f = R_f^{-1}CNOT^{\otimes n}R_f$, where $R_f$ is a reversible implementation of $C_f$.

\kk{expand on above and add any details you deem helpful}























































% \bibliographystyle{alpha}
% \bibliography{main.bib}
% \end{document}

